\documentclass[UTF8, a4paper, oneside, scheme=plain, 12pt]{ctexrep}

\usepackage[T1]{fontenc}
\usepackage{libertinus}
\usepackage{geometry}
\usepackage{float}
\usepackage{titling}
\usepackage{titlesec}
\usepackage{paralist}
\usepackage{footnote}
\usepackage{enumerate}
\usepackage{amsmath, amsthm}
\usepackage{gb4e}
\noautomath
\usepackage{bbm}
\usepackage{textcomp}
\usepackage{soul}
\usepackage{graphicx}
\usepackage{siunitx}
\usepackage[table,xcdraw,svgnames]{xcolor}
\usepackage{tikz}
\usepackage[ruled, vlined, linesnumbered, noend]{algorithm2e}
\usepackage{xr-hyper}
\usepackage[colorlinks, citecolor = purple]{hyperref} % linkcolor=black, anchorcolor=black, citecolor=black, filecolor=black
\usepackage[most]{tcolorbox}
\usepackage{caption}
\usepackage{subcaption}
\usepackage{booktabs}
\usepackage{multirow}
\usepackage[figuresright]{rotating}
\usepackage{acro}
\usepackage[citestyle=authoryear,backend=bibtex,natbib=true,doi=false,isbn=false,url=false]{biblatex}
\addbibresource{references/classical-grammars.bib}
\addbibresource{references/classical-lexicon.bib}
\addbibresource{references/historical-phonology.bib}
\addbibresource{references/general-typology.bib}
\addbibresource{references/other-languages.bib}
\addbibresource{references/language-contact.bib}
\usepackage{prettyref}

\geometry{left=3.18cm,right=3.18cm,top=2.54cm,bottom=2.54cm}
\titlespacing{\paragraph}{0pt}{1pt}{10pt}[20pt]
\setlength{\droptitle}{-5em}

\DeclareMathOperator{\timeorder}{\mathcal{T}}
\DeclareMathOperator{\diag}{diag}
\DeclareMathOperator{\legpoly}{P}
\DeclareMathOperator{\primevalue}{P}
\DeclareMathOperator{\sgn}{sgn}
\newcommand*{\ii}{\mathrm{i}}
\newcommand*{\ee}{\mathrm{e}}
\newcommand*{\const}{\mathrm{const}}
\newcommand*{\suchthat}{\quad \text{s.t.} \quad}
\newcommand*{\argmin}{\arg\min}
\newcommand*{\argmax}{\arg\max}
\newcommand*{\normalorder}[1]{: #1 :}
\newcommand*{\pair}[1]{\langle #1 \rangle}
\newcommand*{\fd}[1]{\mathcal{D} #1}
\newcommand*{\textto}{$\to$}
\newcommand*{\textgt}{$>$ }
\newcommand{\focus}[1]{\textbf{#1}}

\newcommand*{\citesec}[1]{\S~{#1}}
\newcommand*{\citechap}[1]{Ch.~{#1}}
\newcommand*{\citechaps}[1]{Chs.~{#1}}
\newcommand*{\citefig}[1]{Fig.~{#1}}
\newcommand*{\citetable}[1]{Table~{#1}}
\newcommand*{\citepage}[1]{p.~{#1}}
\newcommand*{\citepages}[1]{pp.~{#1}}
\newcommand*{\citefootnote}[1]{fn.~{#1}}

\newrefformat{sec}{\citesec{\ref{#1}}}
\newrefformat{fig}{\citefig{\ref{#1}}}
\newrefformat{tbl}{\citetable{\ref{#1}}}
\newrefformat{chap}{\citechap{\ref{#1}}}
\newrefformat{fn}{\citefootnote{\ref{#1}}}
\newrefformat{box}{Box~\ref{#1}}
\newrefformat{ex}{\ref{#1}}

% color boxes

\tcbuselibrary{skins, breakable, theorems}

\newtcbtheorem[number within=chapter]{infobox}{Box}{
    enhanced,
    boxrule=0pt,
    %colback=blue!5,
    %colframe=blue!5,
    colback=white,
    colframe=white,
    coltitle=blue!60,
    borderline west={4pt}{0pt}{blue!65},
    sharp corners,
    fonttitle=\bfseries, 
    breakable,
    before upper={\parindent15pt\noindent}}{box}
\definecolor{my-orange}{HTML}{F58123}
\newtcbtheorem[number within=chapter, use counter from=infobox]{theorybox}{Box}{
    enhanced,
    boxrule=0pt,
    %colback=orange!5, 
    %colframe=orange!5, 
    colback=white,
    colframe=white,
    coltitle=my-orange!65,
    borderline west={4pt}{0pt}{my-orange!65},
    sharp corners,
    fonttitle=\bfseries, 
    breakable,
    before upper={\parindent15pt\noindent}}{box}
\newtcbtheorem[number within=chapter, use counter from=infobox]{todobox}{Box}{
    enhanced,
    boxrule=0pt,
    colback=red!5,
    colframe=red!5,
    coltitle=red!50,
    borderline west={4pt}{0pt}{red!65},
    sharp corners,
    fonttitle=\bfseries, 
    breakable,
    before upper={\parindent15pt\noindent}}{box}
\newtcbtheorem[number within=chapter, use counter from=infobox]{perspectivebox}{Box}{
    enhanced,
    boxrule=0pt,
    %colback=red!5,
    %colframe=red!5,
    colback=white,
    colframe=white,
    coltitle=red!50,
    borderline west={4pt}{0pt}{red!65},
    sharp corners,
    fonttitle=\bfseries, 
    breakable,
    before upper={\parindent15pt\noindent}}{box}


\AtBeginEnvironment{infobox}{\small}
\AtBeginEnvironment{todobox}{\small}
\AtBeginEnvironment{theorybox}{\small}

\newcommand*{\concept}[1]{\textbf{#1}}
\newcommand*{\term}[1]{\emph{#1}}
\newcommand{\form}[1]{\emph{#1}}
\newcommand{\work}[1]{\textit{#1}}
\newcommand{\species}[1]{\textit{#1}}

\newcommand{\redp}{\textasciitilde}

\DeclareAcronym{blt}{short = BLT, long = Basic Linguistic Theory}
\DeclareAcronym{cgel}{short = CGEL, long = The Cambridge Grammar of the English Language}
\DeclareAcronym{dm}{short = DM, long = Distributed Morphology}
\DeclareAcronym{tag}{long = Tree-adjoining grammar, short = TAG}
\DeclareAcronym{sfp}{long = sentence-final particle, short = \textsc{sfp}}
\DeclareAcronym{np}{long = noun phrase, short = NP}
\DeclareAcronym{vp}{long = verb phrase, short = VP}
\DeclareAcronym{pp}{long = preposition phrase, short = PP}
\DeclareAcronym{cls}{long = classifier, short = CLS}
\DeclareAcronym{dist}{long = distal, short = DIST}
\DeclareAcronym{prox}{long = proximate, short = PROX}
\DeclareAcronym{dem}{long = demonstrative, short = DEM}
\DeclareAcronym{classify}{long = classifier, short = \textsc{cl}}
\DeclareAcronym{dur}{long = durative, short = DUR}
\DeclareAcronym{neg}{long = negative, short = \textsc{neg}}
\DeclareAcronym{cc}{long = copular complement, short = CC}
\DeclareAcronym{cs}{long = copular subject, short = CS}
\DeclareAcronym{tam}{long = {tense, aspect, mood}, short = TAM}
\DeclareAcronym{past}{long = past, short = PST}
\DeclareAcronym{nonpast}{long = non-past, short = NPST}
\DeclareAcronym{present}{long = present, short = PRES}
\DeclareAcronym{progressive}{long = progressive, short = \textsc{poss}}
\DeclareAcronym{perfect}{long = perfect, short = \textsc{perf}}
\DeclareAcronym{passive}{long = passive, short = \textsc{pass}}
\DeclareAcronym{copula}{long = copula, short = COP}
\DeclareAcronym{possessive}{long = possessive, short = \textsc{poss}}

\newcommand{\asis}[1]{\textsc{#1}}
\newcommand{\oneof}[1]{{#1}}
\newcommand*{\homo}[2]{#1$_{\text{#2}}$}

\newcommand{\cgel}{\href{../English/cambridge.pdf}{my notes about CGEL}}
\newcommand{\latin}{\href{../Latin/latin-notes.pdf}{my notes about Latin}}
\newcommand{\alignment}{\href{../alignment/alignment.pdf}{my notes about alignment}}
\newcommand{\exerciseone}{\href{../Exercise/2021-3.pdf}{this exercise}}
\newcommand{\method}{\href{../methodology/glossing.pdf}{this note about my understanding of descriptive grammars}}

\newcommand{\ala}{à la}
\newcommand{\translate}[1]{`#1'}
\newcommand{\vP}{\textit{v}P}
\newcommand*{\category}[1]{\textsc{#1}}
\newcommand*{\specialunit}[1]{$<$\textit{#1}$>$}
\newcommand{\before}{$> \ $}

% Make subsubsection labeled
\setcounter{secnumdepth}{4}
\setcounter{tocdepth}{4}
% reset example counter every chapter (but do not include the chapter number to the label)
\counterwithin{exx}{chapter} 

% Reference formats
\renewcommand*{\nameyeardelim}{\space} % No comma between year and name
\DeclareNameAlias{sortname}{family-given} % Putting the family name before the given name
\DeclareNameAlias{default}{family-given} 
\DeclareFieldFormat{labelnumberwidth}{} % No number label like [12] in the reference list
\setlength{\biblabelsep}{0pt} % No space for these labels

\makeindex

\title{The Early History of Chinese Language}
\author{Jinyuan Wu}

\begin{document}

\maketitle

\chapter{Linguistic evidence}

\section{The prehistory of the Sinitic family}\label{sec:introduction.history}


In this section, we trace the earliest stages of the Sinitic family,
namely the pre-Classical stage, and discuss how it bridges the remote ancestor
and the Classical period.

\subsection{The Trans-Himalayan family}

The Sinitic language family, often collectively referred to as Chinese,
is a branch in the Sino-Tibetan language family, which is also known as the Trans-Himalayan language family.
While the internal structure of the Trans-Himalayan language family
and consequently a reconstruction of Proto-Sino-Tibetan with broad consensus
(at the level of e.g. Proto-Indo-European) remain to be explored,
the shared origin of Chinese, Tibetan, Burmese, and multiple other language families
however is beyond a reasonable doubt:
for example, a list of cognates and a tentative Trans-Himalayan reconstruction can be found in \citet{hill2019historical}.

Having established a ``flat'' relation between Chinese and other languages in the Trans-Himalayan area, 
we may be interested in the position of Chinese within Trans-Himalayan,
and in particular, when it diverged from the rest of the language family.
Note that this is to be decided by comparison of cognates,
and \emph{not} by typological features (\prettyref{sec:intro.sino-tibetan-typology}).
This is inevitably related to the internal structure of the language family.
A consensual phylogeny based on the comparative method is still lacking.
Some scholars take a rather pessimistic view and suggest that
it is probably a better idea to remain agnostic about higher-order subgrouping
and adopt a \emph{Fallen leaves model}, i.e. a thorough enumeration of all genetically certain subgroups without any clear genetic relations stipulated above them.

Two Bayesian phylogenetic studies \citep{sagart2019dated,zhang2019phylogenetic} 
both consider the Sinitic branch to be the earliest branch that diverged from the rest of the language family,
thus supporting the traditional binary division between the Sinitic branch
and the Tibeto-Burmese branch.%
\footnote{
    It should be noted that the huge divergence of contemporary Chinese languages
    from their Tibeto-Burmese cousins 
    cannot be just explained by Neogrammarian sound change laws:
    language contact is an important factor (\prettyref{sec:intro.pre-history.mixing}).
}
They however cannot replace human readable and interpretable comparative analyses,
and currently, the name \term{Tibeto-Burmese}, strictly speaking,
is still not guaranteed to be a node in the phylogenetic tree of the Trans-Himalayan family,
but rather a catch-all term for everything non-Sinitic.

\subsection{On ``Sino-Tibetan typology''}\label{sec:intro.sino-tibetan-typology}

\subsubsection{Purported typological similarities and counterexamples}
Besides cognates, linguists occasionally talk about ``Sino-Tibetan typology''.
Typological features often attributed to Sino-Tibetan languages include the isolating grammar,
tonality, and monosyllabic roots.%
\footnote{
    Modern Sinitic languages are by no means monosyllabic,
    while Classical Chinese appears to be monosyllabic:
    and disyllable words in later stages of Chinese seem to be due to 
    either borrowing or syntactic fossilization of previous analyzable phrases.
    On the other hand, pre-Classical Old Chinese is possibly not strictly monosyllabic
    according to recent reconstructions (\prettyref{sec:intro.reconstruction.minor-syllables}).
}
These criteria was even once used to demarcate the boundary of the Sino-Tibetan language family,
which later proved to be not reasonable.
A refutation on the phonological part can be found in \citet{matisoff1973notes}:
features like tones and monosyllabic roots are not diachronically stable at all. 
As for morphological complexity, just think about how drastically the English inflectional morphology has simplified.
Recent linguistic fieldwork has also revealed that several previously underdocumented Trans-Himalayan languages
completely violate these generalizations.
The Gyalrongic languages, for example, have rather complicated inflectional morphology
and therefore are definitely not monosyllabic,
and Japhug, a Gyalrongic language, lacks tonal contrasts \citep{jacques2021grammar}.

\subsubsection{Typological classification of Proto-Trans-Himalayan}
Given the striking internal typological diversity of Trans-Himalayan languages,
the typological status of Proto-Trans-Himalayan is problematic.
We have reasons to suspect that Proto-Trans-Himalayan is closer to Rgyalrongic in its typological features, actually.
We note that the Rgyalrongic family and the Kiranti family are two distant groups within the Trans-Himalayan family
(which is expected given that they are separated by the Himalayas)
with seemingly cognate personal agreement affixes
which is hard to be explained solely in terms of grammaticalization of pronouns,
and one possible trace of personal agreement is also found in Tibetan
\citep{jacques2012agreement,jacques2010possible,jacques2016tangut}.
Considering the fact that the number of potential cognates between Rgyalrongic and Kiranti is around 150,
and the number of cognates of them with Chinese is of the same order of magnitude,%
\footnote{
    See \citet{shuya2019study} for an explicit historical comparison between Rgyalrongic and Sinitic.
}
it is highly likely that a Rgyalrongic-like personal indexation system originated at a quite early stage,
which possibly is also the ancestor of Chinese
(\citealt{jacques2012agreement}; see \prettyref{box:direct-inverse-history} for more discussions).
Another morphological device that probably has a Proto-Trans-Himalayan origin is the the \form{*s-} prefix,
which is also supported by reconstruction of Old Chinese.
It sometimes is known as the causative prefix,
although this usage is possibly due to grammaticalization of denominal verbs
\citep{jacques2015origin}.

\begin{infobox}{Historical status of the personal agreement affixes}{direct-inverse-history}
    What exactly can be reconstructed to Proto-Sino-Tibetan is not without controversy.
    (a) We may assume that we have a Rgyalrongic-like polypersonal argument indexation system, which is a direct-inverse system, in Proto-Sino-Tibetan.
    (b) Or we can always argues that
    what can be reconstructed to Proto-Sino-Tibetan for certain is just personal affixes,
    as the underlying structure of a direct-inverse system is often not \emph{that}
    different from more common alignments \citep{oxford2017inverse},
    and therefore the direct-inverse system itself possibly cannot be dated back to Proto-Sino-Tibetan
    and only the personal affixes can,
    which are the diachronic materials used to form a Rgyalrongic direct-inverse system in later stages.
    (c) As a step further, it can be argued that the personal agreement affixes
    were probably of pronominal origins, 
    and hence we can argue for a morphologically impoverished Proto-Sino-Tibetan,
    and thus, the great typological distance between Proto-Sino-Tibetan and Chinese is no longer that striking.

    The viewpoints listed above are critically examined in \citet{jacques2012agreement}.
    Although a comprehensive comparison is still lacking,
    \citeauthor{jacques2012agreement} notices that 
    there exist agreement affixes in Rgyalrongic and Kiranti that appear to be cognates and yet cannot be easily explained by grammaticalization of pronouns and hence (c) cannot be completely right,
    thus at least implying certain verbal morphological complexity in a common ancestor.
    Further, by comparative analyses,
    he demonstrates that the verbal conjugation paradigms in Rgyalrongic and in Kiranti are remarkably similar
    (for example, only transitive verbs come with the aforementioned affixes),
    hinting to a biactantial agreement system in the common ancestor, and 
    that Kiranti seems to retain certain residues of an inverse prefix,
    while the evolution of a direct-inverse system, despite theoretically possible,
    remains to be rare in Eurasia.
    The statement (b) above therefore is also not quite feasible, although not completely impossible.
    Anyway, a morphologically rich Proto-Trans-Himalayan is not controversial.
\end{infobox}

\subsection{Materials for reconstructing Pre-Classical Old Chinese}\label{sec:intro-pre-classical}

Middle Chinese phonology was already largely comparable to contemporary Sinitic languages.
Having established that Proto-Trans-Himalayan is morphologically rich,
the evolution from it to today's largely analytic Sinitic languages is rather intriguing.
There has to be a series of transitional stages between an ancestor language with somehow Gyalrong-like typological traits and the modern Sinitic family,
and reconstruction of Old Chinese is essentially
for understanding these stages.

The Old Chinese period is quite long, spanning more than 1,000 years.
Conventionally, the term \term{Old Chinese reconstruction} is mainly about the reconstruction of the language in the \emph{Pre}-Classical period.
Thus, we make use of Pre-Classical texts, including texts from Shang and Zhou dynasties (\prettyref{sec:intro.pre-classical.shang}, \prettyref{sec:intro.pre-classical.post-shang}) and also evidence from neighboring languages,
including early loan words in Old Vietnamese and Japanese.
Finally, modern 闽 Min languages have features that cannot be explained by
the Middle Chinese phonology (see below),
meaning that there exists a stratum in Min languages
which was a sibling, and not a daughter, of Middle Chinese.
This non-Middle Chinese Sinitic branch is often known as Proto-Min.
Finally, Middle Chinese rhyme books and rhyme tables provide a base line
of whatever reconstruction we can think of,
as we expect that the reconstructed Pre-Classical Chinese forms
should explain what we observe in Middle Chinese.

An overview of available materials can be found in \citet[\citechap{2}]{baxter2014old}.
Here we introduce the more conventional sources of Old Chinese,
namely Pre-Classical texts.

\subsubsection{Oracle bones}\label{sec:intro.pre-classical.shang}

The earliest attested Sinitic texts were oracle bone inscriptions from the Shang dynasty.
It should be noted that people during the Shang dynasty \emph{wrote} (and not just \emph{inscribe}) as well:
the counterpart of character 册 exists in oracle bones as well,
which looks a bundle of bamboo and wooden strips.
Still, the only materials that reflects the language of Shang currently available are inscriptions.

The characters in these inscriptions are structured in a way largely comparable to later Chinese characters.
Since the phonetic values of characters in oracle bone inscriptions are not self-evident,
some may even question whether these inscriptions are Sinitic at all:
some may argue that the language of Shang was a (possibly unclassified) SVO language,
and when the Zhou dynasty was established on the ruins of Shang,
its culture was borrowed, including both the writing system and the SVO syntax.
This hypothesis is not quite plausible, because in such a scenario, usually what happens is \emph{lexical} borrowing,
and the introduction of a writing system should be minimal to the structure of the language
\citep[\citepage{88}]{delancey2013origins}.
The introduction of Chinese characters to write Japanese is a good reference point.
Further, phonetic borrowing (假借, \translate{borrowing, i.e. a character originally representing one morpheme to represent another morpheme with a close sound}) in oracle bone inscriptions
neatly fits into reconstructed Old Chinese phonology
(e.g. \citealt[\citepages{115}]{baxter2014old}).


\subsubsection{Post-Shang texts}\label{sec:intro.pre-classical.post-shang}

Oracle bone inscriptions are 
a 20th century archeological re-discovery not known to Classical Chinese authors.
Zhou's replacement of Shang seemed to be accompanied by a great cultural revolution:
the last king of Shang, namely King Zhou of Shang, was known by 
his arrogance, his corruption, and his tyrannic reign.
Yet no one seems to remember the arguably most scary aspect of Shang dynasty: human sacrifice.
To Classical authors, the earliest available texts are 
documents preserved in 《尚书》 (lit. \translate{venerated documents}),
often known as the \work{Book of Documents} in English.
Since these texts are from ancient kings 
whose deeds were romanticized by Confucian scholars,
these texts were highly venerated and yet deemed as
诘屈聱牙 (\translate{twisted, hard to pronounce})
by post-Classical authors.%
\footnote{
    See e.g. 《进学解》 (\work{Analysis of academic advancement}) by Han Yu.
}
They were something that had to be read with commentaries,
the latter written in easier Classical Chinese.
These documents therefore should be regarded as pre-Classical,
although they did contribute sporadic phrases
and grammatical words (e.g. the copula 惟 or the pronoun 厥)
that were occasionally used in Classical Chinese works as a way to polish an article.

A linguistic trait, which corroborates the cultural paradigm shift after the end of the Shang dynasty, 
is that the language of the \work{Book of Documents} and the language of oracle bone inscriptions are clearly different.
The most notable fact on this aspect is that
the aforementioned pronoun 厥 appears frequently in the \work{Documents},
but it appears neither in oracle bone inscriptions nor in Spring and Autumn works. 
Presumably, \work{Book of Documents} contains predominantly early Zhou dynasty texts,
while oracle bones dates back to Shang,
and the differences we are observing reflect dialectal differences between the ruling classes of the two dynasties.

Another fairly early source is 《诗经》 (lit. \translate{poem classics}),
also known as the \work{Book of Odes},
which contains poems dates back to as early as early Zhou.
The \work{Odes}, despite often considered, classical especially in a Confucianism context,
differs greatly from Classical \emph{proses}.
This work mostly concentrate on the grammar of proses.

\subsubsection{Methological issues}\label{sec:intro.prehistory.reconstruction.problems}

None of the texts listed above is alphabetic,
and phonological information is extracted from the texts by observing
(a) information contained in the phonetic components of Chinese characters
(in transmitted texts, and also in excavated texts),
(b) the rhyme patterns in the \work{Book of Odes}.
(a) is done by observing 諧聲 \form{xiéshēng} \translate{harmonic sound} relations,
i.e. semi-homophonous relations observed in Chinese characters without semantic relations that share a component.
Based on inspecting characters with shared components but unrelated meanings in Modern Sinitic languages,
we conclude that what is shared is likely a component representing the \emph{initials} of the words they represent.%
\footnote{
    In modern Sinitic languages, a character usually represents a root,
    but simple philological observation reveals that Old Chinese words were mostly represented by single characters.
}

Due to the uncertainty surrounding the exact phonological information
hidden behind semi-homophonous character series and rhymes in \work{Odes},
as well as the uncertainty of the meaning of certain terms used in 
ancient works on Middle Chinese (e.g. \prettyref{box:middle-chinese-unitary}),
Old Chinese reconstruction is expectedly a chaotic field with lots of controversies.
\citet{baxter2014old} represents an important milestone in this field,
making use of all the materials mentioned above as well as
a rich toolbox of philology, the historical comparative method, and internal reconstruction.
But even it has severe flaws \citep{harbsmeier2016irrefutable}.
Here we list some major methodological problems and not technical errors.

First, currently we do not even have a well-defined target of reconstruction.
Pre-Classical texts show wide variation in syntax and presumably in phonology,
as is shown in \prettyref{sec:intro.pre-classical.shang} and \prettyref{sec:intro.pre-classical.post-shang},
and the Chinese scholar tradition clearly states that 
\work{Book of Odes} consists of poems collected at different times and different locations 
(\prettyref{sec:intro-pre-classical}; \citealt[\citepages{478-480}]{harbsmeier2016irrefutable}).
In the Classical period, texts made it clear that mutually incomprehensible Chinese varieties existed \citep[\citepages{446-447}]{harbsmeier2016irrefutable}.
Thus, a distinction still needs to be made between Old Chinese(s) and Proto-Sinitic,
contrary to the claim of \citet[\citepage{2}]{baxter2014old}.

Second, even when the reconstruction is based on texts known to be from a relatively well-defined \emph{branch} of Old Chinese,
based on e.g. consistent morphosyntactic features,
there is no guarantee that it has a single, unitary phonology
\citep[\citepages{488-490}]{harbsmeier2016irrefutable}.
Note that the same problem occurs for Middle Chinese as well (\prettyref{box:middle-chinese-unitary})

\begin{infobox}{How unitary is Middle Chinese?}{middle-chinese-unitary}
    It should be noted that the problem is not restricted to Old Chinese.
    Strictly speaking, that there was a unitary Middle Chinese
    as reflected by rhyme books and rhyme tables also should not be taken for granted.
    What rhyme books and rhyme tables were intended for could probably never known for certain:
    they could be prescriptive works documenting the mixture of several respected dialects,
    or even an early Neogrammarian-ish reconstruction attempt.
    Besides, we are also not really sure about the phylogenetics of modern dialects:
    if we are to stipulate a Middle Chinese \emph{node} in the Sinitic family tree
    which is the common ancestor of all modern Sinitic dialects,
    then how did it emerge from the messy congeries of dialects during the Middle Chinese \emph{period} \citep[\citepages{477-478}]{harbsmeier2016irrefutable}?
    We are more confident when talking about a unitary Middle Chinese (defined in the rhyme books)
    as the common ancestor of modern Sinitic languages,
    mainly because comparative works do suggest that the former is indeed a good starting point to analyze the phonology of the latter.
    Thus, Middle Chinese in the rhyme books at least is a good proto-language in the Neogrammarian sense.
    Still, Neogrammarian comparative studies based on modern Sinitic languages are urgently needed,
    not mentioning the fact that many are quickly going extinct without proper documentation.
\end{infobox}

Third, what exactly the rhymes and the semi-homophonous relations mean remain not completely clear.
Rhymes can be cultural, containing traditional strata that no longer reflect the actual phonology when the poem was being composed
\citep[\citepages{454-455}]{harbsmeier2016irrefutable}.
The same problem exists for arguments based on shared phonetic components of Chinese characters.
Some, for instance, argue that certain semi-homophonous relations observed
that immediately lead to reconstruction of consonant clusters (\prettyref{sec:intro.prehistory.reconstruction.initial-consonant-cluster})
may simply be due to two pronunciations existing for one character,
which was used to represent two initials because of that.

The severity of this issue however is weaker than the severity of the first two morphological issues, namely the lack of language uniformity in the Pre-Classical Old Chinese period.
First, there is no evidence in the earliest texts that suggests an earlier poem tradition,
and 

Finally, although the validity of certain phonetic series and the related semi-homophonous relations is not clear,
the most prototypical ones like 蓝 vs 监 -- which however are not explainable within the traditional approach -- are well established,
and claiming that one character has multiple pronunciation only pushes the problem to an earlier stage,
as now we have to ask how rather different pronunciations were developed for one character.
Up to now, the only viable explanation is initial consonant clusters.
So at least we are not wrong when saying that consonant clusters existed at a certain stage of Pre-Classical Old Chinese.

\subsection{Outline of Pre-Classical Old Chinese}

Due to all of the issues and difficulties in \prettyref{sec:intro.prehistory.reconstruction.problems}, we will avoid presenting a fine-grained reconstruction of
Old Chinese (Classical or pre-Classical) phonology in this work.
Below, we list some widely accepted features of (pre-Classical) Old Chinese
to give the reader an intuitive and vague feeling of how it feels like.
Understanding how pre-Classical Chinese probably looked like \emph{as a language} and not just as dead texts may provide some insights for understanding the origin of certain features of Classical Chinese, like how valency alternation works.
However, a full survey on the progress of Old Chinese reconstruction is beyond the scope of this work.

\subsubsection{Origin of tones}
In Middle Chinese rhyme books, there are four tones reported:
the \category{level} tone (平聲), the \category{rising} tone (上聲),%
\footnote{
    The term may be confusing for a Modern Standard Mandarin speaker,
    as its reflex in Modern Standard Mandarin is the third tone,
    while the ``raising'' tone in Modern Standard Mandarin is the 陽平 tone,
    which historically was from the \category{level} tone.
}
the \category{departing} tone (去聲), and the \category{entering} tone (入聲).

Let us first focus on the \category{entering} tone.
It can be easily seen that rhymes falling under the category of the \category{entering} tone are much fewer than rhymes pertaining to the other three tones,
with a lot of nasal rhymes missing.
The reflex of roots with \category{entering} tone in modern Sinitic languages, if any,
all end with a stop (consider e.g. Cantonese \form{-p}, \form{-t}, \form{-k}).
These words, when borrowed into Japanese, usually became disyllabic.
These facts can be explained perfectly by assuming that in Middle Chinese,
they ended in \form{-p}, \form{-t}, or \form{-k},
and this seems to be the defining feature of the \category{entering} tone,
which, in Middle Chinese, was strictly speaking not a tone at all.%
\footnote{
    In modern Cantonese there are actually three reflexes of the \category{entering} tone
    and they are real tones,
    but these tones are absent in ancient scholars' analyses of Middle Chinese,
    and are likely recent innovations.
}

Now given that modern Sinitic tones can be easily explained by 
conditional sound changes starting from three tones \citep{sagart1999origin},
the rest three tones in Middle Chinese were real tones.
But we can try to apply the methodology inspired by the analysis above to Old Chinese,
and the conclusion is that there was probably no tone in the Old Chinese period.
Below, I briefly repeat the argumentation in \citet{sagart1999origin}.

The origin of Vietnamese tones has been demonstrated to be final consonants.
That in the oldest layer of Chinese loan words into Vietnamese,
words in the Middle Chinese \category{departing} tone regularly correspond to Vietnamese words in a tone originating from the final consonant \form{-s}
strongly suggests that in Old Chinese,
the origin of the \category{departing} tone is also something like \form{*-s}.
This is supported by textual evidence:
in early Chinese transcriptions of foreign words,
the \category{departing} tone sometimes corresponds to foreign \form{s},
and some early Chinese load words into Korean even retain the \form{s}.
Therefore, the \category{departing} tone is likely a reflex of the Old Chinese word ending consonant \form{*-s}.

In a similar way, the \category{rising} tone is related to an earlier glottal stop
based on evidence from Sino-Vietnamese words.
We also note that \category{rising} tone words were preferred
to transcribe Sanskrit short vowels,
which is consistent with the glottal stop assumption,
as the latter makes the syllable sound short.
Thus, the origin of the \category{rising} tone can be tentatively reconstructed as \form{*-ʔ}.

Now that all other tones have been reduced to syllable final consonants,
the remaining \category{level} tone can be simply understood as the ``default'' case,
i.e. open syllables or syllables ending with nasals.
We have thus shown that Old Chinese lacks tones.

\subsubsection{Affixation}
The fact that Chinese tones probably originated from word-final consonants
and the fact that in reading traditions of Classical texts,
tone change represents subtle modification of the meaning of a root
(known as 四聲别义 \translate{four tones distinguishing between meanings})
immediately suggest that Old Chinese has derivational suffixes.
All derivations that change the tone of a word to the \category{departing} sound,
for instance, can now be explained by a \form{*-s} suffix.
Several prefixes have proposed as well,
based on alternations of the onset,
e.g. multiple \form{*s-} prefixes \citep[\citepages{53-59}]{baxter2014old}.

\subsubsection{Initial consonant clusters}\label{sec:intro.prehistory.reconstruction.initial-consonant-cluster}
The syllable structure of Modern Standard Mandarin and other modern Sinitic languages is relatively simple.
The syllable structure of Middle Chinese was slightly more complicated, allowing more syllable-final consonants (i.e. \category{entering} tone discussed above).
However, all existing materials point to a much more complicated Old Chinese syllable structure.

A popular perspective, although not universally accepted yet, is that consonant clusters were prevalent in Old Chinese.
This can be supported by semi-homophonous series of Chinese characters
(\prettyref{sec:intro.prehistory.reconstruction.problems}).
Consider the semi-homophonous series 监, 槛 (pronounced as \form{jiàn} in the reading convention of Classical texts).
Then consider another series 滥, 蓝.
The fact that they share the same phonetic component 监 strongly suggests that the four characters belong to the \emph{same} series when the characters were first invented.
A reconstructed initial like \form{*kl-} therefore may be proposed for the phonetic value of the phonetic component 监.
A similar ``broken'' series include 禀, 凛 (something like \form{*pl-}).
And now we find a class of consonant cluster initials roughly with the structure of \form{*Cl-},
which actually is reconstructed by \citet[\citepages{50-51}]{baxter2014old} as \form{*Cr-}.

\subsubsection{Final consonant clusters}
Consonant clusters at syllable finals are also proposed.
First, Qing dynasty scholar 段玉裁 Duan Yucai noticed that
in \work{Odes}, rhyme sequences in the \category{level}, \category{rising} and \category{entering} tones
can be isolated (meaning that all rhymes in a poem have the same tone in Middle Chinese),
while the rhyming of \category{departing} tone words appears to be random.
His conclusion was that there was no \category{departing} tone in \work{Odes}.
Now we know that there was probably no tone at all at least in pre-Classical Old Chinese,
and Duan's observation means that
the final consonant \form{*-s} corresponding to the \category{departing} tone
was likely almost always not a part of the root, but something like a suffix.
But if a word with the \category{entering} tone (which means it ends with \form{*-p}, \form{*-t}, or \form{*-k})
rhymed with a word that ends with \form{*-s},
then the latter should also contain \form{*-p}, \form{*-t}, or \form{*-k} near its end:
and hence the \form{*-ps}, \form{*-ks}, \form{*-ts} final consonant clusters
\citep[\citepage{183}]{hill2019historical}.

In Middle Chinese, there are four rhymes (祭泰夬废) in the \category{departing} tone
that do not have counterparts in the \category{level} and \category{rising} tones,
and words with these rhymes having semi-homophonous relations with \category{entering} words,
based on \form{xiéshēng} relations seen in their characters.
This means these rhymes probably ended with \form{*-ps}, \form{*-ts}, \form{*-ks} in Old Chinese,
and have both \category{entering} and \category{departing} features.
Consonant clusters appear here again.

We also note that there are \category{rising} tone words with nasal finals in Middle Chinese:
上 itself is an instance.
This suggests a  consonant cluster roughly with the form of \form{*-ŋʔ} at the end of the words,
which probably should be \form{*-ʔŋ} to make it looks typologically feasible.
We note that the form \form{*-ʔŋ} can even explain the origin of the \category{rising} tone,
as it makes the pronunciation \emph{weak} in the middle,
which naturally is close to a pronunciation that is \emph{low} in the middle,
and hence the modern \category{rising} tone.

\subsubsection{Minor syllables}\label{sec:intro.reconstruction.minor-syllables}
The format of the \work{Odes} strongly suggests that
the language(s) behind the poems were monosyllabic.
Still, we have evidence for residues of minor syllables before the main syllables,
known as \emph{loose} pre-initials in \citet{baxter2014old};
their \emph{tight} counterparts are ordinary consonant clusters.
In \form{Odes}, some rather strange forms, like 帝命不時,
contain 無 that would otherwise be interpreted as a negator,
but the interpretations of them contain no negative meanings
according to commentaries.
Not all of them can be naturally explained as rhetorical questions.
We can also find records like 聿……吳謂之不律 in 说文,
further suggesting the existence of loose pre-initials
\citep[\citepages{122-1235}]{hill2019historical}.
The reconstruction of these pre-initials is still controversial,
but their existence at least in some historical stages and areas seems to be beyond reasonable doubts.
The syllable structure of pre-Classical Old Chinese therefore should be (C\form{ə})CVC.


\subsubsection{Word order}
Texts during the Old Chinese period do \emph{not} have a rather flexible word order, as is shown in the rest of this work.
The evolution towards a largely analytic syntax therefore had to already have started at the Old Chinese stage.





\subsection{Language contact}\label{sec:intro.pre-history.mixing}

\subsubsection{Neighboring languages}

Certain records in the Classical period indicate the existence of non-Sinitic populations in southern kingdoms.
For example, the 越人歌 \form{Yuèréngē} \translate{\work{Song of Boatmen of Yue}},
a song of boatmen preserved in \work{Garden of Stories} that was both transcribed in Chinese characters (in a sort of Old Chinese pronunciation) and translated,
was proposed to be composed in a Kra-Dai language \citep{zhengzhang1991decipherment}.
What we are discussing above happened long \emph{before} the Classical period,
but the relevant populations did not come out of no where
and should be there as well when Proto-Sinitic was being formed.

Major non-Sino-Tibetan language families in touch with 
include Kra-Dai, Hmong–Mien, and Vietic.
Among them, 

What is particularly intriguing is that around 20\% of words in \citet{mccraw2010lexical} do not have a clear origin.
These words are often collectively referred to as having \form{Dongyi} 東夷 \translate{eastern babarians} origins in modern Chinese scholarship,
because these words have to come from language contact,
and the most likely candidate of the origin is these ``barbarians'' i.e. non-Sinitic populations in the west of Shang and Zhou.
Therefore, it seems that a ghost population once existed between the northern Sinitic population and the southern Austronesian population,
speaking a language that heavily influenced Proto-Sinitic.
No known surviving descendant of Dongyi is known. 


\subsubsection{The makeup of the Old Chinese lexicon}
Less than a half of Old Chinese words have clear etymological Tibeto-Burman counterparts,
and the rest of the lexicon comes from Austro-Asiatic and Tai contributions,
with the former being the majority
\citep{mccraw2010lexical}.
Whether this means Old Chinese was more heavily influenced by Austro-Asiatic 
than other Trans-Himalayan languages is unclear,
as a comparison between Austro-Asiatic loan words in Trans-Himalayan is currently lacking.
Also worth noticing is that although in later periods, Chinese (at different stages of development) was the \emph{origin} of borrowed terms,
in the prehistoric era, Chinese was frequently the \emph{recipient}, not the \emph{donor}.

\subsubsection{Syntax of Old Chinese and language contact}
A natural idea is that the syntax of Old Chinese -- and hence the syntax of Classical Chinese -- 
was also influenced by non-Trans-Himalayan languages.
The SVO word order, the prepositions, and clause-initial subordinators all stand in contrast to Tibeto-Burman languages.
On the other hand, Old Chinese has sentence-final particles, which seem to be more consistent with a SOV order,
and occasionally we observe relics of a SOV order as in 不我知也;
these features are not surprising within the Trans-Himalayan family.

How sentence-final particles were grammaticalized is not always clear:
in oracle bone inscriptions there are already some SFPs;
see http://ir.lib.scu.edu.tw:8080/bitstream/987654321/11005/1/103SCU00045056-001.pdf:
句尾亦偶有少數的語氣詞

We note that the all of the Southeast Asian groups with SVO syntax seem to
have this order as far back as we can trace,
and therefore in the aspect of word order (and probably other aspects),
Chinese is the \emph{recipient}, not the \emph{donor} \citep[\citepage{82-85}]{delancey2013origins}.

\subsection{Discussion: how Chinese as we know it was formed?}

To summarize what has been revealed in reconstruction works,
absence of tones, words with minor syllables,
and various affixation mechanisms were characteristics of Chinese 
at a certain stage \citep[\citepage{13}]{sagart1999roots}.
In the earliest Chinese texts, however, we already see indications of
emergence of stereotypical Sinitic features:
the syntax was already quite rigid,
the number of available affixes was limited and their productivity was likely weakened,
and the syllable structure was under simplification (TODO: ref).
The factors pushing Chinese on its path to its appearance today are of interest.

\subsubsection{Relation with Chinese characters}

It is sometimes argued that the loss of morphology is due to the Chinese characters,
which do not record inflection well.
This claim is linguistically untenable,
first because most people before modernity were illiterate,
and second because it is quite straightforward to dedicate certain characters to represent inflection
(which is indeed how the \form{kana} system used to write Japanese developed).
A reasonable, but not consensually accepted hypothesis is that it was due to language contact,
which is discussed in \prettyref{sec:intro.pre-history.mixing}.

\subsubsection{Area typology}
Certain languages in the Trans-Himalayan area that do not have clear genetic relations with the Trans-Himalayan family, on the other hand,
do have the purported Trans-Himalayan typological features.
Kra-Dai languages have these typological features as well \citep[\citepage{434}]{sidwell2021languages},
and yet the most likely relatives of them are the Austronesian family,
which do \emph{not} have these typological features at all \citep{sagart2004higher,ostapirat2005notes}.

The evolution of tonal distinctions happens in Kra-Dai as well \citep{sagart2019model}.


\end{document}